% Tabela consolidada de resultados: processed_v3 vs dataset_indt
% Detector siamese de erros de layout (cabeca sobre DINOv2 ViT-S/14 congelado).
% Numeros: artifacts/reports/EXPERIMENT_RESULTS.json (processed_v3) e
%          artifacts/indt/reports/EXPERIMENT_RESULTS.json (dataset_indt).
% Pacotes necessarios no preambulo:
%   \usepackage{booktabs}
%   \usepackage{threeparttable}
%   \usepackage{multirow}   % (opcional)

% ============================ TABELA I (principal) ============================
\begin{table*}[t]
\centering
\caption{Consolidated comparison of the two evaluation datasets for the layout-error
detector (siamese head over a frozen DINOv2 ViT-S/14). Held-out (test) split, threshold
calibrated on validation. \textbf{Bold} marks the better value \emph{only} among rows that
are cross-dataset comparable.}
\label{tab:dataset-comparison}
\begin{threeparttable}
\begin{tabular}{l l c c}
\toprule
\textbf{Group} & \textbf{Metric} & \textbf{\texttt{processed\_v3}} & \textbf{\texttt{dataset\_indt}} \\
\midrule
\multirow{3}{*}{\textit{Protocol}}
 & Eval.\ clean source            & real (same device) & synthetic (app, reflow)\tnote{a} \\
 & Test images (clean / error)    & 108 (41 / 67)      & 184 (24 / 160) \\
 & Resolution-only baseline AUROC & 1.000\tnote{b}     & 0.497 \\
\midrule
\multicolumn{4}{l}{\textit{Confound-free detection} -- errors injected into clean, same resolution (\textbf{primary, comparable})} \\
 & AUROC & \textbf{0.721} & 0.671 \\
 & AP    & \textbf{0.903} & \textbf{0.903} \\
\midrule
\multicolumn{4}{l}{\textit{Error categorisation} -- Stage~2, prototype decider} \\
 & Coarse macro-F1 (2 super-cls.) [95\% CI] & \textbf{0.626} [0.51, 0.74] & 0.488 [0.40, 0.58] \\
 & Fine macro-F1 (4 classes)                & \textbf{0.336}             & 0.303 \\
\midrule
\multicolumn{4}{l}{\textit{Operating point and confound diagnostics} -- \textbf{not cross-dataset comparable}\tnote{c}} \\
 & Accuracy                          & 0.583 & 0.391 \\
 & Recall (error)                    & 0.582 & 0.300 \\
 & Specificity                       & 0.585 & 1.000\tnote{d} \\
 & Global gate AUROC (fusion)        & 0.596 & 0.774 \\
 & Falsifiability: error / resolution& 0.596 / 0.596\tnote{e} & 0.774 / 0.290 \\
\bottomrule
\end{tabular}
\begin{tablenotes}[flush]
\footnotesize
\item[a] \texttt{dataset\_indt} has \emph{no} real device clean screens; the clean class is
synthetic (legitimate-reflow/benign variants of held-out app screens). App parents are
partitioned per app into disjoint train/val/test sets to avoid resubstitution.
\item[b] In \texttt{processed\_v3} every real clean screen shares one resolution
(2076$\times$2152), so a trivial ``resolution $\neq$ canonical $\Rightarrow$ error'' rule
already reaches AUROC~1.0; the global gate metrics are therefore confound-dominated.
\item[c] These rows depend on the clean-negative \emph{domain}, which differs between datasets
(real same-device vs.\ synthetic app), so they must not be compared head-to-head.
\item[d] Computed on 24 \emph{synthetic} clean screens (app domain), not on real device clean;
it does not certify a zero real-world false-alarm rate. The low recall (0.300) reflects the
synthetic$\rightarrow$real transfer gap.
\item[e] AUROC of the model score predicting \emph{error} vs.\ predicting the resolution
confound; a near-zero gap (\texttt{processed\_v3}) means the score tracks resolution as much as
error. The large gap in \texttt{dataset\_indt} is partly an artefact of the synthetic clean
removing the resolution split (see note~b).
\end{tablenotes}
\end{threeparttable}
\end{table*}

% ============================ TABELA II (por classe) ============================
\begin{table}[t]
\centering
\caption{Per-class results on the held-out split: Stage~1 detection (AUROC vs.\ clean) and
Stage~2 categorisation (macro-F1). \textbf{Bold} = better per row.}
\label{tab:per-class-comparison}
\begin{threeparttable}
\begin{tabular}{l c c c c}
\toprule
 & \multicolumn{2}{c}{\textbf{Detection AUROC}\tnote{a}} & \multicolumn{2}{c}{\textbf{Class.\ F1}} \\
\cmidrule(lr){2-3}\cmidrule(lr){4-5}
\textbf{Class} & \texttt{proc.\_v3} & \texttt{indt} & \texttt{proc.\_v3} & \texttt{indt} \\
\midrule
\texttt{black\_bars}        & 0.728 & \textbf{0.869} & \textbf{0.611} & 0.563 \\
\texttt{disordered\_layout} & 0.627 & \textbf{0.756} & 0.000          & \textbf{0.227} \\
\texttt{empty\_space}       & 0.517 & \textbf{0.696} & 0.222          & \textbf{0.327} \\
\texttt{overlay}            & 0.531 & \textbf{0.742} & \textbf{0.510} & 0.095 \\
\bottomrule
\end{tabular}
\begin{tablenotes}[flush]
\footnotesize
\item[a] Detection AUROC is category-vs-clean and is \emph{confounded} (each class has its own
resolution profile, and the clean negative differs in domain between datasets); read it as
indicative, not as a head-to-head ranking. \texttt{black\_bars} is the strongest class on both
datasets; \texttt{overlay} the weakest to classify.
\end{tablenotes}
\end{threeparttable}
\end{table}
